\documentclass[14pt]{extarticle}
\usepackage{natbib}

\usepackage{xcolor}
\usepackage{amsmath}
\newcommand{\tuple}[1]{ \langle #1 \rangle }
\usepackage{times}
\usepackage{ltablex}
\usepackage{tasks}
\usepackage{threeparttable}
\usepackage{booktabs}
\usepackage{xcolor}
\usepackage{tikz}
\usepackage{tcolorbox}
\usepackage{amsmath, amsfonts, amssymb, amsthm}
\usepackage{natbib}
\usepackage{iftex}
\ifPDFTeX
  \usepackage[T1]{fontenc}
  \usepackage{lmodern}
\else
  \usepackage{fontspec}
  \usepackage{luatexja}
\fi
\usepackage[mathscr]{euscript}

%%%%%% Template
\usepackage{hyperref}
\hypersetup{colorlinks=true,allcolors=blue}

\usepackage{vmargin}
\setpapersize{USletter}
\setmarginsrb{1.0in}{1.0in}{1.0in}{0.6in}{0pt}{0pt}{0pt}{0.4in}

\sloppy

\usepackage{parskip}
\setlength{\parindent}{0cm}

\allowdisplaybreaks

\newcommand{\orange}[1]{\textcolor{orange}{#1}}
\newcommand{\blue}[1]{\textcolor{blue}{#1}}
\newcommand{\red}[1]{\textcolor{red}{#1}}

\def\qqquad{\quad\qquad}
\def\qqqquad{\qquad\qquad}

%%%%%%%%%%%%%%%%%%%%%%%%%%%%%%%%%%%%%%%%%%%%%%%%%%%%%%%%%%%%%%%%%%%%%%%%%%%%%%%%
% fill-in-blank question style, found in https://tex.stackexchange.com/a/505089

\usepackage{ifthen}
\usepackage{tocloft}
\usepackage{exercise}

% Set the Show Answers Boolean
\newboolean{showAns}
\setboolean{showAns}{false}
\newcommand{\showAns}{\setboolean{showAns}{true}}

% The length of the Answer line
\newlength{\answerlength}
\newcommand{\anslen}[1]{\settowidth{\answerlength}{#1}}

% ans command that indicates space for an answer or shows the answer in red
\newcommand{\ans}[1]{\settowidth{\answerlength}{\hspace{2ex}#1\hspace{2ex}}%
    \ifthenelse{\boolean{showAns}}%
        {\textcolor{red}{\underline{\hspace{2ex}#1\hspace{2ex}}}}%
        {\underline{\hspace{\answerlength}}}}%

\newcommand{\details}[1]{\settowidth{\answerlength}{#1}%
    \ifthenelse{\boolean{showAns}}%
        {\begin{tcolorbox}[colback=blue!5!white,colframe=blue!75!black,title=Solution] #1 \end{tcolorbox}}%
        {}}%

% Formatting how multiple choices Questions are formated.
\settasks{label=(\Alph*), label-width=30pt}

% Some commands for the Exercise Question package
\renewcommand{\QuestionNB}{\Large\protect\textcircled{\small\bfseries\arabic{Question}}\ }
\renewcommand{\ExerciseHeader}{} %no header
\renewcommand{\QuestionBefore}{3ex} %Space above each Q
\setlength{\QuestionIndent}{8pt} % Indent after Q number

% To create the list of answers with tocloft...
\newcommand{\listanswername}{Answers}
\newlistof[Question]{answer}{Answers}{\listanswername}

% Creates a TOC for Answers
\newcounter{prevQ}
\newcommand{\answer}[1]{\refstepcounter{answer}%
\ans{#1}%
\ifnum\theQuestion=\theprevQ%
        \addcontentsline{Answers}{answer}{\protect\numberline{}#1}%
        \else%
        \addcontentsline{Answers}{answer}{\protect\numberline{\theQuestion}#1}%
        \setcounter{prevQ}{\value{Question}}%
        \fi%
        }%

\sloppy

\newcommand{\HRule}{\rule{\linewidth}{0.5mm}}
\newcommand{\Hrule}{\rule{\linewidth}{0.3mm}}

%tocloft formatting listofanswers
\renewcommand{\cftAnswerstitlefont}{\bfseries\large}
\renewcommand{\cftanswerdotsep}{\cftnodots}
\cftpagenumbersoff{answer}
\addtolength{\cftanswernumwidth}{10pt}

\makeatletter
\renewcommand{\@maketitle}{%
  \parindent=0pt%
  \centering
  {\Large \bfseries\textsc{\@title}} \\
  \vspace{5pt}
  {\large \textit{\@author}} \\
  \HRule \\
  \vspace{1em}
}
\makeatother

\ifPDFTeX
  % pdfLaTeX: Latin Modern (no fontspec)
\else
  % LuaLaTeX: use Crimson Pro if available; otherwise fall back to Latin Modern
  \IfFontExistsTF{Crimson Pro Light}{%
    \setmainfont{Crimson Pro Light}[
      ItalicFont={* Italic},
      BoldFont={Crimson Pro Medium},
      BoldItalicFont={Crimson Pro Medium Italic}]
    \setsansfont{Crimson Pro Light}[
      ItalicFont={* Italic},
      BoldFont={Crimson Pro Medium},
      BoldItalicFont={Crimson Pro Medium Italic}]
  }{%
    \setmainfont{Latin Modern Roman}
    \setsansfont{Latin Modern Sans}
  }%
\fi

%%%%%%%%%%%%%%%%%%%%%%%%%%%%%%%%%%%%%%%%%%%%%%%%%%%%%%%%%%%%%%%%%%%%%%%%%%%%%%%%
% Uncomment next line to show answers in red:
% \showAns
% Uncomment next line to print a list of answers:
% \listofanswer

\title{Problem Set 1}
\author{Hui-Jun Chen}

\begin{document}
\maketitle

\begin{Exercise}

\subsection*{Lecture 1: Motivation --- Inflation, the price level, and nominal anchors}

\Question The \textbf{price level} $P_t$ is best interpreted as \answer{B}
\begin{tasks}(1)
\task the growth rate of prices between $t$ and $t+1$
\task the value of money in terms of goods
\task the nominal interest rate set by the central bank
\task the real interest rate that clears the goods market
\end{tasks}

\Question \textbf{Inflation} $\pi_t$ is \answer{A}
\begin{tasks}(1)
\task the rate of change of the price level
\task the level of real GDP
\task the nominal money supply
\task the real wage divided by productivity
\end{tasks}

\Question A \textbf{nominal anchor} is needed because otherwise \answer{C}
\begin{tasks}(1)
\task output is not uniquely determined in the long run
\task the labor market cannot clear
\task the price level can drift without a force pinning down its level
\task the government cannot issue nominal debt
\end{tasks}

\Question In the 2021--2022 U.S. episode, one headline puzzle motivating the course is \answer{D}
\begin{tasks}(1)
\task inflation fell immediately once deficits rose
\task real GDP was constant while inflation changed
\task policy rates rose sharply before inflation rose
\task inflation rose substantially while policy rates initially stayed near zero
\end{tasks}

\Question The fiscal-theory/debt-valuation equation conceptually says the real value of nominal government liabilities equals \answer{B}
\begin{tasks}(1)
\task the current money growth rate
\task the present value of expected future primary surpluses
\task the present value of expected future GDP growth
\task the current-period tax revenue only
\end{tasks}

\Question A key reason \emph{``deficits cause inflation''} is not a complete theory is that \answer{C}
\begin{tasks}(1)
\task deficits always reduce output
\task deficits can only be financed by money printing
\task what matters is the expected \emph{path} of future fiscal backing, not just current deficits
\task deficits mechanically determine the real interest rate
\end{tasks}

\Question In the debt-valuation view, unexpected inflation can be triggered by \answer{A}
\begin{tasks}(1)
\task news that lowers expected future primary surpluses (backing)
\task higher current productivity with no other changes
\task a fall in money demand with fixed $P$
\task an increase in labor supply holding everything else fixed
\end{tasks}

\Question Under a regime where the public expects future taxes/spending to adjust to stabilize debt, a temporary deficit is most consistent with \answer{D}
\begin{tasks}(1)
\task large and permanent inflation regardless of policy
\task an immediate collapse in output
\task a permanent fall in the price level
\task relatively muted inflation pressure from the deficit alone
\end{tasks}

\Question A clean way to summarize the course motivation is: theories of inflation must explain \answer{B}
\begin{tasks}(1)
\task only the long-run growth rate of output
\task timing and comovement of $\pi$, $i$, output, and policy during episodes
\task only cross-sectional variation in firm prices
\task only the money multiplier
\end{tasks}

\Question In Cochrane-style framing, disagreements about inflation often boil down to disagreements about \answer{C}
\begin{tasks}(1)
\task whether labor supply is elastic
\task whether investment has adjustment costs
\task the policy regime and the nominal anchor
\task whether production is Cobb--Douglas
\end{tasks}

\subsection*{Lecture 2: IS--LM --- money market, goods market, and aggregate demand}

\Question The money-demand function $M^d/P = L(Y,i)$ typically satisfies \answer{A}
\begin{tasks}(1)
\task $\partial L/\partial Y>0$ and $\partial L/\partial i<0$
\task $\partial L/\partial Y<0$ and $\partial L/\partial i>0$
\task $\partial L/\partial Y=0$ and $\partial L/\partial i<0$
\task $\partial L/\partial Y>0$ and $\partial L/\partial i=0$
\end{tasks}

\Question If output $Y$ rises (holding $M$ and $P$ fixed), the LM curve implies the equilibrium interest rate \answer{B}
\begin{tasks}(1)
\task falls because money demand falls
\task rises because money demand rises
\task stays constant because money supply is fixed
\task becomes indeterminate because goods market clears first
\end{tasks}

\Question Holding nominal money supply $M$ fixed, an increase in the price level $P$ \answer{C}
\begin{tasks}(1)
\task raises real balances $M/P$ and shifts LM right
\task does not affect LM because $M$ is nominal
\task lowers real balances $M/P$ and shifts LM left
\task shifts IS right through higher government spending
\end{tasks}

\Question A monetary expansion ($M\uparrow$) with $P$ fixed shifts the LM curve \answer{D}
\begin{tasks}(1)
\task up/left
\task not at all
\task up/right
\task down/right
\end{tasks}

\Question In standard IS logic, a fall in the interest rate increases equilibrium output because \answer{A}
\begin{tasks}(1)
\task investment (and other interest-sensitive spending) rises
\task money demand rises
\task productivity rises
\task labor supply falls
\end{tasks}

\Question A rightward shift of IS can be caused by \answer{B}
\begin{tasks}(1)
\task an increase in money demand holding everything else fixed
\task an increase in government purchases $G$
\task an increase in the price level $P$ holding $M$ fixed
\task a fall in productivity $A$
\end{tasks}

\Question IS--LM equilibrium is the point where \answer{C}
\begin{tasks}(1)
\task the labor market and goods market clear
\task only the goods market clears
\task both the goods market and money market clear
\task money demand equals investment demand
\end{tasks}

\Question The \textbf{aggregate demand} (AD) curve slopes downward because \answer{A}
\begin{tasks}(1)
\task $P\downarrow \Rightarrow M/P\uparrow \Rightarrow$ LM shifts right $\Rightarrow Y\uparrow$
\task $P\downarrow \Rightarrow$ wages rise $\Rightarrow Y\uparrow$
\task $P\uparrow \Rightarrow$ productivity rises $\Rightarrow Y\uparrow$
\task $P\downarrow \Rightarrow$ taxes rise $\Rightarrow Y\uparrow$
\end{tasks}

\Question In the IS--LM-to-AD derivation, which variable is treated as changing to trace out AD? \answer{D}
\begin{tasks}(1)
\task $Y$
\task $i$
\task $G$
\task $P$
\end{tasks}

\Question In the short run of the textbook IS--LM story, it is common to treat the price level $P$ as \answer{B}
\begin{tasks}(1)
\task fully flexible within the period
\task predetermined/sticky so IS--LM pins down $Y$ given $P$
\task irrelevant because money is neutral
\task chosen by households directly
\end{tasks}

\subsection*{Lecture 3: AD--AS --- demand vs supply shocks and short-run vs medium-run}

\Question The FE (full-employment) line in $(Y,P)$ space is typically \answer{C}
\begin{tasks}(1)
\task downward sloping because higher $P$ raises real balances
\task upward sloping because higher $P$ raises labor supply
\task vertical because full-employment output is determined by real factors
\task horizontal because prices are fixed
\end{tasks}

\Question A rightward shift in AD (holding AS fixed) typically leads to \answer{A}
\begin{tasks}(1)
\task $Y\uparrow$ and $P\uparrow$ in the short run
\task $Y\uparrow$ and $P\downarrow$ in the short run
\task $Y\downarrow$ and $P\uparrow$ in the short run
\task $Y\downarrow$ and $P\downarrow$ in the short run
\end{tasks}

\Question A negative productivity shock is best represented in AD--AS as \answer{D}
\begin{tasks}(1)
\task AD shifting right
\task AD shifting left
\task SRAS shifting right
\task SRAS shifting left (and FE shifting left)
\end{tasks}

\Question A supply shock that shifts AS left tends to generate \answer{B}
\begin{tasks}(1)
\task $Y\uparrow$ and $P\downarrow$
\task $Y\downarrow$ and $P\uparrow$
\task $Y\uparrow$ and $P\uparrow$
\task $Y\downarrow$ and $P\downarrow$
\end{tasks}

\Question The comovement diagnostic says: demand shocks typically make $(Y,P)$ move \answer{A}
\begin{tasks}(1)
\task in the same direction
\task in opposite directions
\task not at all
\task randomly
\end{tasks}

\Question In the medium-run adjustment story, if $Y>Y^{FE}$ after an AD expansion, then price adjustment typically implies \answer{C}
\begin{tasks}(1)
\task the price level falls until $Y$ rises further
\task $Y$ stays permanently above $Y^{FE}$
\task the price level rises, reducing real balances and moving $Y$ back toward $Y^{FE}$
\task the money supply must fall one-for-one with output
\end{tasks}

\Question A fiscal expansion $G\uparrow$ shifts AD right. In a model with a vertical FE line, the long-run effect on output is \answer{D}
\begin{tasks}(1)
\task permanently higher output
\task permanently lower output
\task output becomes indeterminate
\task no permanent change in output; mainly a higher price level
\end{tasks}

\Question In AD--AS, an outward shift of AD combined with an inward shift of AS is most consistent with \answer{B}
\begin{tasks}(1)
\task disinflation and a boom
\task higher inflation with ambiguous output (depends on magnitudes)
\task lower inflation with lower output
\task unchanged inflation and higher output
\end{tasks}

\Question The main reason Lecture 3 follows Lecture 2 is that IS--LM gives AD, but \answer{A}
\begin{tasks}(1)
\task inflation requires a supply/price-adjustment side (AS)
\task the labor market disappears in IS--LM
\task money demand is irrelevant for policy
\task GDP accounting fails
\end{tasks}

\Question In a simple AD--AS framework, ``stagflation'' refers to \answer{C}
\begin{tasks}(1)
\task low inflation and high output
\task low inflation and low output
\task high inflation and low output
\task high inflation and high output
\end{tasks}

\subsection*{Lecture 4: Inflation dynamics --- Fisher, Phillips curves, policy rules, and Cochrane-style distinctions}
\Question The Fisher equation can be summarized as \answer{B}
\begin{tasks}(1)
\task $i_t \approx r_t - E_t\pi_{t+1}$
\task $i_t \approx r_t + E_t\pi_{t+1}$
\task $r_t \approx i_t + E_t\pi_{t+1}$
\task $E_t\pi_{t+1} \approx r_t - i_t$ and is always zero
\end{tasks}

\Question Under adaptive expectations $\pi_t^{e}=\pi_{t-1}$, the ex-ante real rate is \answer{C}
\begin{tasks}(1)
\task $r_t = i_t + \pi_{t-1}$
\task $r_t = i_t - \pi_{t+1}$
\task $r_t = i_t - \pi_{t-1}$
\task $r_t = i_t$ always
\end{tasks}

\Question A backward-looking Phillips curve with a cost-push shock is closest to \answer{D}
\begin{tasks}(1)
\task $\pi_t = \beta E_t\pi_{t+1} + \kappa x_t$
\task $\pi_t = \pi_{t-1} - \kappa x_t$
\task $\pi_t = E_t\pi_{t+1} + u_t$
\task $\pi_t = \pi_{t-1} + \kappa x_t + u_t$
\end{tasks}

\Question Rearranging $\pi_t = \pi_{t-1} + \kappa x_t + u_t$ implies \answer{A}
\begin{tasks}(1)
\task $x_t = (\pi_t-\pi_{t-1}-u_t)/\kappa$
\task $x_t = (\pi_{t-1}-\pi_t-u_t)/\kappa$
\task $x_t = (\pi_t+\pi_{t-1}+u_t)/\kappa$
\task $x_t = \kappa(\pi_t-\pi_{t-1}) + u_t$
\end{tasks}

\Question Suppose policy follows a Taylor rule $i_t = \bar i + \phi_\pi \tilde\pi_t + \phi_x x_t$ and $\pi_t^e=\pi_{t-1}$.
Then the real-rate gap can be written as \answer{B}
\begin{tasks}(1)
\task $r_t-r^n = \phi_\pi \tilde\pi_{t-1} + \phi_x x_t - \tilde\pi_{t}$
\task $r_t-r^n = \phi_\pi \tilde\pi_{t} + \phi_x x_t - \tilde\pi_{t-1}$
\task $r_t-r^n = \phi_\pi \tilde\pi_{t} - \phi_x x_t + \tilde\pi_{t-1}$
\task $r_t-r^n = \tilde\pi_{t} - \tilde\pi_{t-1}$ only
\end{tasks}

\Question With IS in gap form $x_t=-\sigma(r_t-r^n)+\varepsilon_t^d$, substituting the expression in the previous question and solving for $x_t$ gives \answer{D}
\begin{tasks}(1)
\task $x_t = -\sigma\phi_\pi\tilde\pi_t + \varepsilon_t^d$
\task $x_t = \sigma\phi_\pi\tilde\pi_t - \sigma\tilde\pi_{t-1} + \varepsilon_t^d$
\task $x_t = \dfrac{-\sigma\phi_\pi\tilde\pi_{t-1}+\sigma\tilde\pi_t+\varepsilon_t^d}{1+\sigma\phi_x}$
\task $x_t = \dfrac{-\sigma\phi_\pi\tilde\pi_t+\sigma\tilde\pi_{t-1}+\varepsilon_t^d}{1+\sigma\phi_x}$
\end{tasks}

\Question Combining IS + backward-looking Phillips curve + Taylor rule yields an AR(1) for inflation deviations $\tilde\pi_t=a\tilde\pi_{t-1}+\cdots$.
In this system, a sufficient condition for mean reversion ($a<1$) is \answer{C}
\begin{tasks}(1)
\task $\phi_\pi<1$
\task $\phi_\pi=1$
\task $\phi_\pi>1$
\task $\phi_x=0$ always
\end{tasks}

\Question A positive demand shock $\varepsilon_t^d>0$ raises inflation in the Old-Keynesian system primarily because \answer{B}
\begin{tasks}(1)
\task it directly increases $\pi_t$ even if $x_t=0$
\task it raises $x_t$ (output gap), which feeds into inflation via $\kappa x_t$
\task it lowers the policy rate mechanically through the Fisher equation
\task it shifts full-employment output $Y^{FE}$ immediately
\end{tasks}

\Question A positive cost-push shock $u_t>0$ means \answer{A}
\begin{tasks}(1)
\task inflation rises even if $x_t=0$; offsetting it requires $x_t<0$ (slack)
\task inflation rises only if the policy rate falls
\task inflation must fall because firms become more productive
\task inflation is unchanged because $u_t$ cancels in equilibrium
\end{tasks}

\Question Cochrane emphasizes ``distinguishing theories'' because different models imply different predictions about \answer{D}
\begin{tasks}(1)
\task the accounting identity $Y=C+I+G$
\task whether $x_t$ is defined as $Y_t-Y^{FE}$
\task whether CPI or PCE should be used
\task inflation under interest-rate pegs/Taylor rules and the role of fiscal backing/news
\end{tasks}

\end{Exercise}

\end{document}
